\documentclass[12pt]{article}
\usepackage[utf8]{inputenc}
\usepackage[T1]{fontenc}
\usepackage{hyperref}
\title{Toba-Maskoy Language Resources for K-12 Curriculum Development}
\author{Compiled by Claude Sonnet 5 (claude-sonnet-5) for Andr\'as Kornai}
\date{}
\begin{document}
\maketitle

\section{Introduction}
Toba-Maskoy (ISO 639-3 tmf, Glottocode toba1268), also called Enenlhet, is an endangered Mascoian (Enlhet--Enenlhet) language of the Paraguayan Chaco. Sources give about 1,700 speakers. The searches returned background pages but no dictionary, grammar or school materials.

\section{Primary Educational Resources}
None were found.

\section{Linguistic and Reference Materials}
\begin{itemize}
\item Toba-Maskoy language (Wikipedia): \url{https://en.wikipedia.org/wiki/Toba-Maskoy_language}
\item Mascoian languages (Wikipedia mirror): \url{https://en.wikipedia.com/wiki/Mascoyan}
\item Toba-Maskoy numerals data (Kalisch). The page points to Unruh, Kalisch and Romero (2003) on numeral syntax, which I did not locate: \url{https://lingweb.eva.mpg.de/channumerals/Toba-Maskoy.htm}
\item A journal article found in the results (Sínteses, Unicamp). I did not verify its content. \url{https://revistas.iel.unicamp.br/index.php/sinteses/article/view/1234/918}
\end{itemize}

\section{Religious and Community Resources}
None were found.

\section{Digital and Audio Resources}
Catalogue entry in the Endangered Languages Project: \url{https://endangeredlanguages.com/lang/2141?hl=fr}

\section{Curriculum Development Context}
Little is available online. The Enlhet--Enenlhet family has related varieties, and any curriculum work should be done with the communities and with researchers who have documented Enenlhet. Those sources need further searching outside this report, and they may exist but were not found here.
\end{document}